\section{预量子化、极化与半形式：有限维相空间上的三级台阶}
\label{sec:prequant}

\subsection{这一章要解决什么}

上一章结束时，我们有了一个有限维、紧、带奇点的辛空间 $\cM$（$\Sigma$ 上平坦联络模规范变换），
以及它上面的辛形式
\begin{equation}
\omega=\frac{k}{4\pi}\int_\Sigma\Tr\bigl(\delta A\wedge\delta A\bigr)
\qquad\text{（\eqref{eq:ch3-omega}）}.
\end{equation}
但“一个辛流形”本身给不出 Hilbert 空间。几何量子化要把“辛流形”变成“态空间”，中间恰好有三级台阶：

\begin{enumerate}[label=\textbf{台阶 \arabic*.},leftmargin=3.2em]
\item \textbf{预量子化}：构作一个线丛 $\cL$ 与联络，使曲率正比于 $\omega$。这一步有\emph{拓扑障碍}：
若 $[\omega]$ 不代表整上同调类，线丛根本不存在。
\item \textbf{极化}：截面空间太大（等价于“同时依赖位置与动量”），必须选一半方向，只在那一半上取
协变常数截面。
\item \textbf{半形式修正}：量子化要求线丛的“平方根张量积”，于是有效次数/能级发生平移。
\end{enumerate}

本文 §3.1 只用了一段话跨过这三级台阶，且\textbf{没有提到第一级有障碍}。本章把这三级补出来；
补的过程会发现，本文其实是用“行列式线丛”这把钥匙把第一级台阶\emph{绕过}了，而不是不需要它。

\subsection{原文说了什么}

原文相关的只有四处，逐条列出（锚点与强度见表 \ref{tab:ch4-statements}）。

\begin{table}[htbp]
\centering
\small
\setlength{\tabcolsep}{4pt}
\begin{tabular}{@{}p{0.05\textwidth}p{0.63\textwidth}p{0.14\textwidth}p{0.11\textwidth}@{}}
\toprule
编号 & 原文陈述 & 锚点 & 强度 \\
\midrule
T1 & 选一个复结构 $J$，$\cM$ 成为复流形，且是某类全纯丛的模空间 & \src{368--369} & 陈述（引 [33]）\\
T2 & 行列式线丛 $\cL$ 定义在 $\cM_J$ 上；$\cL$ 的第一陈类即 $k=1$ 的辛形式 & \src{369} & 陈述 \\
T3 & 一般 $k$ 取 $\cL$ 的 $k$ 次幂；量子态空间是全纯截面空间 & \src{369} & 陈述 \\
T4 & $\cM$ 紧时该截面空间有限维 & \src{369} & 陈述（与上章 S6 同一论证）\\
\bottomrule
\end{tabular}
\caption{本文关于预量子化的四条陈述。原文没有出现“整性条件”“极化”“半形式”这三个词。}
\label{tab:ch4-statements}
\end{table}

原文的核心句子（英文原句）：

\begin{quote}\small
``\ldots once one picks a complex structure $J$ on $\Sigma$, the moduli space $\cM$ of flat
connections can be given a new interpretation --- it is the moduli space of a certain family of
holomorphic vector bundles.'' \src{369}\\[2pt]
``\ldots the space of global holomorphic sections of $\cL^{\otimes k}$.'' \src{369}
\end{quote}
中间那一步——为什么可以从“行列式线丛”直接得到“量子态空间”——原文只用一段说明：$\cL$ 来自
$\bar\partial$-算子族的行列式，Dirac 行列式给它一个自然度规，取 $k$ 次幂即得预量子化线丛
\src{369}，并由此得到
\begin{equation}
\cH_\Sigma^{(J)}=H^0\bigl(\cM_J,\cL^{\otimes k}\bigr).
\label{eq:ch4-hilbert}
\end{equation}
本章的任务就是把“$\Rightarrow$”填满。

\subsection{原文没说的细节}




\medskip
\noindent\textbf{G4.1}\quad \textbf{预量子化的整性条件，以及本文为什么不检查它。}
标准陈述（Weil--Kostant）：辛流形 $(M,\omega)$ 存在预量子化线丛 $\iff$
$\frac{1}{2\pi\hbar}[\omega]\in H^2(M;\mathbb Z)$。原文没有这条。读者要么误以为 CS 情形不需要整性，
要么误以为作者漏了。正确说法是：本文用行列式线丛\emph{构造性地}给出了线丛，而它的第一陈类天生
是整的，所以整性在这个具体问题里是“免费”的——但这是\emph{结构}带来的，不是理论不需要它。



\medskip
\noindent\textbf{G4.2}\quad \textbf{行列式线丛的定义、Quillen 度规、曲率公式。}
原文只写“determinant line bundle of the $\bar\partial$-operator”与“the Dirac determinant gives a
natural metric” \src{369}，没有定义 $\det H^0\otimes(\det H^1)^{-1}$，没有写 Quillen 度规的公式，
也没有写“曲率 = $\frac{\ii}{2\pi}\omega$”。这三条是整章的技术核心，缺了就无法判断 $\cL$ 与 $\omega$
是否匹配（差 $2\pi$ 还是差 $2\pi k$，全在这里）。



\medskip
\noindent\textbf{G4.3}\quad \textbf{极化的定义与三种类型。}
“选 $J$”在几何量子化的语言里就是“取 Kähler 极化”。原文只说“picks a complex structure $J$”
\src{368}，没有定义极化（拉格朗日 + 可积 + 维数恒定），也没有说为什么紧相空间上不能取 $q,p$ 分离的
实极化。



\medskip
\noindent\textbf{G4.4}\quad \textbf{“量子态 = 沿极化协变常数的截面”这一句本身。}
原文直接说“全纯截面” \src{369}。但“全纯”只是极化条件 $\nabla^{0,1}s=0$ 在平凡化后的读法，
中间要用到 Kähler 条件（$F^{0,2}=0$）才能把 $(0,1)$-部分的联络当成 $\bar\partial$ 算子。



\medskip
\noindent\textbf{G4.5}\quad \textbf{半形式修正与 $k\to k+h^\vee$。}
本文 §3 完全没提半形式，§4 却处处用 $q=\exp\frac{2\pi\ii}{N+k}$（(4.21)，\src{382}）。
这个平移从哪来？在几何量子化里它是“张量积 $K^{1/2}$”；在本文里它以另一种面貌出现在 §2 的 $\eta$
不变式里（(2.16)(2.17)，\src{359}）。两条路的关系要讲清，而且要说清\emph{本文没有做半形式这一条}。



\medskip
\noindent\textbf{G4.6}\quad \textbf{Narasimhan--Seshadri 定理。}
原文用一句话把“平坦联络模空间”等同于“稳定全纯丛模空间” \src{368--369}（引 [33]）。
这是本书最关键的等同之一（也是第 \ref{sec:blocks} 节“共形块 = 态空间”的基础），
但原文没有陈述定理内容。


\subsection{补出的推导}

\noindent\textbf{D-4.1（Weil--Kostant 陈述，G4.1）}\quad\emph{起点}：经典的 Weil--Kostant 陈述（$[\omega]/(2\pi\hbar)$ 整）；\emph{终点}：四步造出 $(\cL,\nabla)$，$\nabla^2=-\frac{\ii}{\hbar}\omega$。
设 $(M,\omega)$ 是辛流形，$\hbar>0$。几何量子化的第一步要一个 Hermite 线丛 $(\cL,h,\nabla)$
使得
\begin{equation}
F_\nabla=-\frac{\ii}{\hbar}\,\omega ,
\label{eq:ch4-preq}
\end{equation}
其中 $F_\nabla=\nabla^2$ 是曲率 2-形式 [李思 p.302，定义 5.4.2]。这是唯一的相容性要求：
它保证“沿 Hamiltonian 向量场的协变化”与“泊松括号”在截面上对易子级别一致，即
$\bigl[\nabla_{X_f},\nabla_{X_g}\bigr]-\nabla_{[X_f,X_g]}=-\ii\hbar\,\omega(X_f,X_g)$
[李思 p.302]。由此定义预量子化算子 $Q(f)=\ii\hbar\nabla_{X_f}+f$，即“乘以函数 + 沿流协变导数”。

\begin{derivation}[D-4.1：从整性条件造出预量子化线丛（四步）]
\textbf{起点}：辛流形 $(M,\omega)$ 与常数 $\hbar$；\textbf{终点}：一个 Hermite 线丛 $(\cL,\nabla)$ 满足
\eqref{eq:ch4-preq}。
\begin{enumerate}[label=\textbf{步 \arabic*.},leftmargin=2.8em]
\item \textbf{局部 1-形式。} $\dd\omega=0$ 且 $M$ 可缩的局部片上，$\omega=\dd\theta_\alpha$。取好覆盖
$\{U_\alpha\}$，在重叠处 $\theta_\alpha-\theta_\beta=\dd f_{\alpha\beta}$。
\item \textbf{周期量化。} 整性条件 $\frac{1}{2\pi\hbar}[\omega]\in H^2(M;\mathbb Z)$ 等价于：可选取
$f_{\alpha\beta}$ 使 $\frac{1}{\hbar}f_{\alpha\beta}$ 在 $\mathbb Z$ 上取模的“平移”满足
$\oint_\gamma\omega\in 2\pi\hbar\mathbb Z$（$\gamma$ 为任意闭 2-周期）。这是“$\theta$ 的平移只能按
$2\pi\hbar$ 的整数倍”这一句话的全局版本。
\item \textbf{拼线丛。} 令转移函数 $g_{\alpha\beta}=\exp\bigl(\frac{\ii}{\hbar}f_{\alpha\beta}\bigr)$。
Cech 闭合性给出 $g_{\alpha\beta}g_{\beta\gamma}g_{\gamma\alpha}=1$，于是 $\{g_{\alpha\beta}\}$ 定义线丛
$\cL$；局部联络 $\nabla_\alpha=\dd-\frac{\ii}{\hbar}\theta_\alpha$ 在重叠处相容
（$\nabla_\alpha=g_{\alpha\beta}\nabla_\beta g_{\alpha\beta}^{-1}$）。
\item \textbf{验曲率。} $\nabla^2=-\frac{\ii}{\hbar}\dd\theta=-\frac{\ii}{\hbar}\omega$，即 \eqref{eq:ch4-preq}。
于是 $F_\nabla$ 的 de Rham 类 $=\frac{1}{2\pi\hbar}[\omega]$，而它必须落在 $H^2(M;\mathbb Z)$ 里——这就是
必要性。充分性（给定整性类能真的拼出线丛）由 Weil 与 Kostant 给出，此处按陈述级引用
[李思 p.303--304]。
\end{enumerate}
\end{derivation}

这个方程有解的必要条件是：$\omega$ 的 de Rham 类必须落在整格上。因为 $F_\nabla$ 是线丛联络的曲率，
$\bigl[\frac{\ii}{2\pi}F_\nabla\bigr]=c_1(\cL)\in H^2(M;\mathbb Z)$，而由 \eqref{eq:ch4-preq}
左边正是 $\frac{1}{2\pi\hbar}[\omega]$。于是

\begin{keypoint}
\textbf{整性条件（Weil--Kostant）}：存在满足 \eqref{eq:ch4-preq} 的预量子化线丛，当且仅当
\begin{equation}
\frac{1}{2\pi\hbar}\,[\omega]\in H^2(M;\mathbb Z).
\end{equation}
必要性是陈类整性；充分性由 Weil 与 Kostant 给出（$\omega$ 的整性 + $M$ 的辛形式局部可写
$\sum\dd p_i\wedge\dd q_i$ 的拼接）。李思的教材陈述了条件并说“反过来也对”，但\textbf{没有给出
充分性的证明} [李思 p.303--304]。本章据此把整性记为“陈述级引用”，不冒充定理证明。
\end{keypoint}

在 CS 情形上，$M=\cM$，$\hbar$ 已被吸收进 $k$（见 D-4.2），条件变成“$[\omega]$ 是 $2\pi$ 的整数倍”；
而 $\omega$ 的系数是 $\frac{k}{4\pi}$，所以 $k\in\mathbb Z$ 就是这个条件——这与 §1 里由大规范变换得到的
$k\in\mathbb Z$ \src{354} 是同一件事的两种来源。本文选择了后者（作用量的单值性），完全没有走前者。

\noindent\textbf{D-4.2（为什么本文不需要检查整性：行列式线丛的构造性绕过，G4.1--G4.2）}\quad\emph{起点}：原文只有“取行列式线丛的 $k$ 次幂”一句；\emph{终点}：$c_1(\cL)$ 由行列式构造自动整，故无需检查整性。
行列式线丛的定义如下。设 $\pi:\mathcal{E}\to\cM$ 是模空间上的全纯向量丛族，纤维 $E_m=H^\bullet(\Sigma;E_m)$
的上同调随 $m$ 变维数。定义
\begin{equation}
\cL=\det H^0(\Sigma;E)\otimes\bigl(\det H^1(\Sigma;E)\bigr)^{-1},
\label{eq:ch4-det}
\end{equation}
它是 $\cM$ 上的全纯线丛；在 $k=1$ 时，原文断言其第一陈类就是辛形式 \src{369}。用整性的语言：
\begin{equation}
[\omega]_{k=1}=2\pi\,c_1(\cL)\in H^2(\cM;2\pi\mathbb Z),
\end{equation}
即 $[\omega]$ 的整性\emph{自动}成立。这就是“构造性绕过”的确切含义 \ours：
不是整性条件不成立，而是构造线丛的同时把陈类算出来了，条件被满足而无需检验。

对一般的 $k$，令 $\cL_k=\cL^{\otimes k}$，则 $F_{\cL_k}=kF_{\cL}=-\ii k\omega_{k=1}=-\ii\omega$（按
$\omega=\frac{k}{4\pi}\int$ 的归一化），与 \eqref{eq:ch4-preq} 对照即 $\hbar=1$。所以
\begin{equation}
\boxed{\;\text{本文的全套归一化等价于取 }\hbar=1,\quad \cL^{\otimes k}\text{ 即预量子化线丛。}\;}
\end{equation}
李思书里的“$L^{\otimes k}$、$F=kF_0=-\ii(2\pi k)\omega$、认同 $\hbar=\frac{1}{2\pi k}$” [李思 p.308，
例 5.4.11] 是同一个构造换了 $\hbar$ 的位置；两种写法给出相同的物理（$k\to\infty$ 即 $\hbar\to0$ 的
半经典极限）。这一条对读本文极重要：本文从不出现 $\hbar$，全部半经典信息在 $k$ 上。

\noindent\textbf{D-4.3（Quillen 度规与曲率，G4.2）}\quad\emph{起点}：$\cL$ 上的 Hermite 度规；\emph{终点}：Quillen 度规及其曲率 $\frac{\ii}{2\pi}\omega$（陈述级）。
在 \eqref{eq:ch4-det} 上取 Quillen 度规：对 $s\in\det H^0$，
\begin{equation}
\|s\|_Q^2=\|s\|_{L^2}^2\cdot\det{}'\bigl(\bar\partial^*\bar\partial\bigr)^{1/2}
\quad\text{（乘一个与 $s$ 无关的 $\zeta$-正则化因子）}.
\end{equation}
它给出的 Chern 形式就是 $\omega$ 的那一半：
\begin{equation}
c_1(\cL,\|\cdot\|_Q)=\Bigl[\frac{\ii}{2\pi}\omega\Bigr].
\label{eq:ch4-curv}
\end{equation}
\eqref{eq:ch4-curv} 是陈述级引用：它属于 Quillen（1985）与 Freed（1987）的工作，在带反常的族上
需要 Bismut--Freed 的推广。本书没有逐句核对这些文献，也不在此假装给出证明；读者只需知道：
\textbf{$\cL$ 与 $\omega$ 的匹配不是猜的，而是 $\bar\partial$ 族的行列式的标准性质}。

原文把它压缩成一句“the Dirac determinant gives a natural metric” \src{369}，并且明确说对
$SU(N)$ 以外的群要作修正（脚注 9 提到 $E_8$ 的情形，其中某个线丛的 30 次幂才等于 $\cL$）——这一点
说明“行列式线丛给预量子化”在一般群上不是免费的，本文自己也知道。我们只核到该脚注的首句与
$E_8$ 的例子，细节未核（见本章边界）。

\noindent\textbf{D-4.4（极化：三种选择，G4.3）}\quad\emph{起点}：相空间上的复结构 $J$；\emph{终点}：极化的三种类型，以及量子化 $=$ 沿 $P$ 的协变常数截面。
极化是 $TM$ 的复子丛 $P$ 满足：(a) 拉格朗日：$\omega|_P=0$，$\dim_\mathbb C P=\frac12\dim_\mathbb R M$；
(b) 可积：$[P,P]\subset P$；(c) 维数恒定 [李思 p.305，定义 5.4.6]。量子态定义为沿 $P$ 协变常数的截面：
\begin{equation}
\nabla_X s=0\qquad\forall X\in\Gamma(P).
\label{eq:ch4-pol}
\end{equation}
三种典型选择：
\begin{enumerate}[label=(\alph*),leftmargin=2.2em]
\item \textbf{实极化} $P=\mathrm{span}\{\partial_{p_i}\}$：条件变成“$s$ 与 $p$ 无关”，
即 $s=s(q)$——这就是 Schrödinger 表象。
\item \textbf{Kähler 极化}：$P=T^{0,1}M$（复结构 $J$ 的负全纯切丛）。\eqref{eq:ch4-pol} 变成
“$s$ 在 $J$-全纯坐标下全纯”。
\item \textbf{混合/不极化}：不做。
\end{enumerate}
原文取的是 (b)：“picks a complex structure $J$” \src{368} 就是取 $P=T^{0,1}\cM_J$。

\noindent\textbf{D-4.5（Kähler 极化给出全纯截面，G4.4）}\quad\emph{起点}：K\"ahler 极化 $P=T^{0,1}M$；\emph{终点}：量子态空间 $=H^0(\cM_J,\cL^{\otimes k})$ 的全纯截面，有限维。
这一步要验证“协变常数 = 全纯”。因为 $P=T^{0,1}$ 是 $TM$ 的一半，$\nabla$ 的 $(0,1)$ 部分
$\nabla^{0,1}$ 定义了一个 $\bar\partial$-算子；$s$ 是 $\nabla^{0,1}$-常数当且仅当在局部全纯平凡化下
$\bar\partial s=0$。关键是：$\nabla^{0,1}$ 平方为零——
\begin{equation}
\bigl(\nabla^{0,1}\bigr)^2=F^{0,2}=0 ,
\end{equation}
这\emph{正是} Kähler 条件（$(M,J,\omega)$ 的 Kähler 形式闭 $\Rightarrow$ 曲率的 $(0,2)$ 部分为零）。
于是量子态空间就是
\begin{equation}
\cH^{(J)}=H^0\bigl(\cM_J,\cL^{\otimes k}\bigr),
\end{equation}
即 \eqref{eq:ch4-hilbert}。李思把这条链讲得比原文细：先证相容性
$[\nabla_X,\nabla_Y]-\nabla_{[X,Y]}=-\ii\hbar\omega(X,Y)$，再取极化，再取截面 [李思 p.305--307]；
例 5.4.9 给出 Kähler 情形即全纯截面、例 5.4.10 给出 $P=T^{0,1}M\Rightarrow H=H^0(M,L)$
（且紧 Kähler 时有限维）[李思 p.306--307]。

\begin{keypoint}
\textbf{为什么不能取 $q,p$ 分离的实极化。} 在 $\mathbb R^{2n}$ 上 (a) 型极化给出 $L^2(\mathbb R^n)$，
很好用。但在\emph{紧}相空间上，$P=T^{0,1}$ 之外的选择会遇到障碍：拉格朗日叶若是闭的，
\eqref{eq:ch4-pol} 的截面在一整片叶上只能差一个常数，而叶是紧的 $\Rightarrow$ 只有满足
Bohr--Sommerfeld 条件 $\int_\gamma\theta\in2\pi\hbar\mathbb Z$（$\theta$ 为联络 1-形式）的叶才可有
非零截面；一般叶上截面为零。Kähler 极化没有这个麻烦，因为 $T^{0,1}$ 的叶就是点，条件变成“局部全纯”，
而全纯截面在紧复流形上总是有限维的 \ours。原文没有解释这一点，直接跳到 Kähler 极化
\src{368--369}，原因就在这里。
\end{keypoint}

\noindent\textbf{D-4.6（Narasimhan--Seshadri，G4.6）}\quad\emph{起点}：平坦联络与全纯丛两套语言；\emph{终点}：Narasimhan--Seshadri 的稳定全纯丛 $\leftrightarrow$ 平坦联络。
定理（Narasimhan--Seshadri，1965）：设 $\Sigma$ 是紧 Riemann 面，$G$ 紧，$G_\mathbb C$ 为其复化。
次数为零的稳定全纯 $G_\mathbb C$-丛与 $\Sigma$ 上平坦 $G$-联络模规范变换一一对应：
\begin{equation}
\cM=\Hom(\pi_1\Sigma,G)/G\;\cong\;\cM_{\mathrm{stable}}(\Sigma,G_\mathbb C).
\end{equation}
对应关系的方向：给定平坦联络，$\bar\partial_A$ 给出全纯结构，平坦性 $\Rightarrow$ 该丛稳定；
反方向用 Hermite--Einstein 度量解 $\bar\partial$ 方程。对 $G=SU(N)$ 这就是“稳定 rank $N$、次数 $0$
的全纯丛模空间” \src{368--369}。它给出：
\begin{itemize}[leftmargin=1.6em]
\item $\cM$ 上的复结构 $J$（从而可以谈 Kähler 极化）；
\item $\cM$ 是射影代数簇（因此是紧的、有丰富的全纯线丛）；
\item 上章 D-3.7 的“奇点”在代数侧对应“非稳定/可约丛”的边界。
\end{itemize}
原文只给结论并引 [33] \src{368--369}；定理本身的页码未逐字核（见边界）。

\noindent\textbf{D-4.7（半形式修正与 $k\to k+h^\vee$，G4.5）}\quad\emph{起点}：$K_{G/T}=\cL_{-2\rho}$、$K_{\cM_J}\simeq\cL^{-2h^\vee}$；\emph{终点}：半形式使有效次数平移 $\mp h^\vee$，与 $k+c_2(G)/2$ 对齐。
几何量子化的严格版本要求“半形式”：量子态不是 $\cL^{\otimes k}$ 的截面，而是
$\cL^{\otimes k}\otimes K^{1/2}$ 的截面，$K$ 是极化方向的典范丛（Kähler 极化下 $K=K_{\cM_J}$，
即全纯切丛行列式的对偶）。理由是量子化要与“两次量子化”相容（metaplectic 修正）：辛形式的一半
“$\sqrt{\dd p\wedge\dd q}$”也要张量进来。

在旗流形 $G/T$ 上这条公式还能算到底：$K_{G/T}=\cL_{-2\rho}$（$\rho$ 为半正根和，$\cL_\lambda$ 为
权 $\lambda$ 的齐次线丛），故
\begin{equation}
K_{G/T}^{1/2}=\cL_{-\rho}.
\label{eq:ch4-half}
\end{equation}
于是对 $G/T$ 做“$k$ 级 + 半形式”的量子化，等价于做“$k-\rho$ 级”的量子化。这个计算对 Borel--Weil--Bott
部分（第 \ref{sec:wilson} 节）是常用的，也解释了为什么表示的最高权会自动带上 $\rho$ 的平移。

\begin{derivation}[D-4.7：半形式把 $\lambda$ 平移 $\rho$——逐步论证]
\textbf{起点}：K\"ahler 极化下的量子化 $H^0(M,\cL^{\otimes k})$；\textbf{终点}：半形式修正使有效次数
平移 $h^\vee$，与 §2 的 $k\to k+c_2/2$ 对齐。
\begin{enumerate}[label=\textbf{步 \arabic*.},leftmargin=2.8em]
\item \textbf{量子化对象换成“线丛 $\otimes$ 半形式”。} 严格量子化取
$\cL^{\otimes k}\otimes K_P^{1/2}$，$K_P$ 为极化方向典范丛；K\"ahler 极化 $P=T^{0,1}M$ 时
$K_P\simeq K_M$（全纯 $n$-形式丛）。
\item \textbf{旗流形上算到底。} $M=G/T$：$K_{G/T}$ 对应最高权 $-2\rho$ 的齐次线丛，即
$K_{G/T}=\cL_{-2\rho}$，故 $K_{G/T}^{1/2}=\cL_{-\rho}$（这就是 \eqref{eq:ch4-half}）。于是
\begin{equation}
H^0(G/T,\cL_\lambda\otimes K^{1/2})
=H^0(G/T,\cL_{\lambda-\rho}).
\end{equation}
\item \textbf{用 Borel--Weil--Bott 读出平移。} BWB 说 $H^0(G/T,\cL_\mu)=V_\mu$（$\mu$ 支配整）或 $0$
（否则）。所以“带半形式、在 $\lambda$ 处量子化”$=$“不带半形式、在 $\lambda-\rho$ 处量子化”。
反过来：要得到表示 $V_\lambda$，辛形式必须取 $\lambda+\rho$ 那一类。所谓“$\rho$ 平移”就是指这个
从 $\lambda$ 到 $\lambda+\rho$ 的移动。
\item \textbf{搬到模空间。} 把上一步的平移语言搬到 $\cM_J$：线丛 $\cL$ 的第一陈类是 $[\omega]$ 的
$1/(2\pi)$ 倍（\eqref{eq:ch4-curv}），而 $K_{\cM_J}$ 按 \eqref{eq:ch4-kmod} 是 $\cL$ 的
$-2h^\vee$ 次幂。于是 $\cL^{\otimes k}\otimes K^{1/2}$ 的有效次数是 $k\mp h^\vee$，符号由
$c_1(\cL)=\pm[\omega]/(2\pi)$ 的取向约定决定。这一步是\emph{类比}，不是定理：我们把 $G/T$ 上严格成立的
$\rho$ 平移按“同一套几何量子化语言”迁移到 $\cM_J$ 上，标为[我们]。
\item \textbf{与单圈计算对齐。} 本文 §2 的 $\eta$ 不变式给出的是 $k\to k+c_2(G)/2$ (2.16)(2.17)
\src{359}；而 $c_2(G)=2h^\vee$ 正是本文的归一化（$SU(N)$：$c_2=2N$，$h^\vee=N$）。两条数字因此相符：
半形式给出的是同一个 $h^\vee$ 平移，只是符号与归一化要靠约定固定。
\end{enumerate}
\end{derivation}

\textbf{但本文没有做这一步}，而且这里有一个符号陷阱。若把同样的关系搬到 $\cM_J$ 上，用模空间上
常见的陈述
\begin{equation}
K_{\cM_J}\simeq\cL^{\otimes(-2h^\vee)}\qquad\text{（陈述级，未核原始文献）},
\label{eq:ch4-kmod}
\end{equation}
则“$\cL^{\otimes k}\otimes K^{1/2}$”等于 $\cL^{\otimes(k-h^\vee)}$，看起来是把 $k$ 往\emph{下}平移；
而物理上要的是 $k\to k+h^\vee$。差号的来源是 $c_1(\cL)$ 与 $[\omega]$ 的符号约定（$\cL$ 是正线丛
还是负线丛，取决于把 $\omega$ 写成 $2\pi c_1(\cL)$ 还是 $-2\pi c_1$），本文从未固定这一约定。

\begin{keypoint}
\textbf{正确的做法是：不要从半形式去“推”$k+h^\vee$，而是记住两条独立证据。}
\begin{itemize}[leftmargin=1.4em]
\item 证据一（本文自己的）：§2 的单圈/$\eta$ 不变式计算给出耦合的平移
$k\to k+c_2(G)/2$ (2.16)(2.17) \src{359}。对 $SU(N)$，$c_2(SU(N))=2N$，故平移 $N$。
\item 证据二（几何量子化）：半形式给出的是“$\cL^{\otimes k}\otimes K^{1/2}$”，其效果是一个
$h^\vee$ 级的平移（符号依约定）。
\end{itemize}
两条数字相符，因为本文在 §2 用的 $c_2$ 归一化正满足 $c_2(G)=2h^\vee$（这就是 (2.17) 里
$c_2(G)/2=h^\vee$ 的来历）。所以“有效能级是 $k+h^\vee$”是一个\emph{被两条独立计算支持的结论}，
而不是本文给出的一条定理；本文的正式推导走的是 $\eta$ 不变式，半形式只是另一种语言。
\end{keypoint}

\noindent\textbf{D-4.8（$k$ 的角色与半经典极限）}\quad\emph{起点}：$k$ 的三个角色（整数、线丛次数、维数参数）；\emph{终点}：$q=\exp\frac{2\pi\ii}{k+h^\vee}$ 与 $k\to\infty$ 的半经典意义。
把上面的内容合起来：相空间固定（$\cM$，与 $k$ 无关），量子化随 $k$ 变。$k$ 同时扮演三个角色：
\begin{enumerate}[label=(\roman*),leftmargin=2.2em]
\item 拓扑量子化参数：$k\in\mathbb Z$ 来自大规范变换下 $e^{\ii L}$ 的单值性 \src{354}；
\item 预量子化线丛的次数：$\cL^{\otimes k}$，$k\to\infty$ 即 $\hbar\to0$（D-4.2）；
\item 态空间维数的控制参数：$\dim\cM$ 不变，$\dim\cH^{(J)}$ 随 $k$ 增长（$\sim k^{g}$ 级别，
见 \eqref{eq:ch4-hilbert} 的 Riemann--Roch 计数）[李思 p.308--309]，且要求 $k\ge$ 某阈值才有
非零态空间（上章 G3.5 提到的“小 $k$ 维数可小于经典值”）。
\end{enumerate}
这三重身份不矛盾，但混起来读就会出错——例如以为“$k$ 大”意味着“相空间大”，恰恰相反：相空间不变，
变的是线丛的次数。

把三重身份合到一个公式里，就是二维侧的量子参数（原文 (4.21) \src{382}，一般 $G$ 写为 $k+h^\vee$）：
\begin{equation}
q=\exp\frac{2\pi\ii}{k+h^\vee},
\qquad\text{其中 } SU(N):\ h^\vee=N,\quad SU(2):\ h^\vee=2 .
\label{eq:q}
\end{equation}
本章的贡献只是解释这个平移的来源：\emph{不是} $k$ 变成 $k+h^\vee$ 的相空间变了，而是预量子化线丛
被半形式（或等价地，被 §2 的 $\eta$ 不变式）改了次数。

\subsection{与其他来源的对照}

\begin{table}[htbp]
\centering
\small
\setlength{\tabcolsep}{4pt}
\begin{tabular}{@{}p{0.14\textwidth}p{0.24\textwidth}p{0.26\textwidth}p{0.26\textwidth}@{}}
\toprule
主题 & 本文 & 李思 5.4 & 河野 2.5 / Moore 3.4 \\
\midrule
整性条件 & 不出现；用行列式线丛绕过 \src{369} & 陈述 $\frac{1}{2\pi\hbar}\int_\Sigma\omega\in\ZZ$，
必要性 + “反过来也对”，充分性无证明 [李思 p.303--304] & 河野用 $[\omega]\in H^2(M;\ZZ)$ 与
$\nabla=\dd-2\pi\sqrt{-1}\alpha$ 的约定 [河野 p.viii--ix]\\
预量子化线丛 & 行列式线丛 $\cL$，$k$ 次幂 \src{369} & 定义 5.4.2：$F_\nabla=-\frac{\ii}{\hbar}\omega$
[李思 p.302] & Moore 3.4 用“平坦模空间的量子化”重述\\
极化 & “picks $J$” \src{368} & 定义 5.4.6（拉格朗日+可积+维数恒定）[李思 p.305] & 河野 2.5 半形式语言\\
全纯截面 & $H^0(\cM_J,\cL^{\otimes k})$ \src{369} & 例 5.4.9（Kähler）、例 5.4.10（$H^0(M,L)$ 有限维）
[李思 p.306--307] & Moore 用 $\mu^{-1}(0)/G$ 的量子化\\
半形式 & 未做 & 未讲 half-form 机制（只讲 $\hbar=1/(2\pi k)$ 的换算，例 5.4.11 [李思 p.308]） &
河野 2.5 有半形式与 level shift 的讨论（页码按 PDF−14，未逐字核）\\
\bottomrule
\end{tabular}
\caption{预量子化三级台阶的三方对照。最后一行的“未逐字核”是诚实的标注：\\
本书只核到河野 §2.5 的位置（书 p.$\sim$84--86），没有逐句核对其半形式公式。}
\label{tab:ch4-sources}
\end{table}

\subsection{陷阱与约定}

\begin{keypoint}
\textbf{五个约定坑。}
\begin{enumerate}[leftmargin=1.6em]
\item \textbf{$\hbar$ 在哪。} 本文 $\hbar=1$（D-4.2），李思保留 $\hbar$，河野把 $2\pi$ 显式写进联络。
三者可以互相换算，但\textbf{不要在同一个公式里混用}。
\item \textbf{$\cL^{\otimes k}$ 与 level。} 本文的“level $k$”是\emph{带引号}的简写（原文
``\ldots `level' $k$ \ldots'' \src{370}），其真值是仿射 level $k+h^\vee$。§3 里用 $k$ 无害，
因为那里只做有限维相空间的量子化；一旦和流代数/共形块对齐（第 \ref{sec:blocks} 节），
必须换成 $k+h^\vee$。
\item \textbf{$\cL$ 是正还是负。} 由 \eqref{eq:ch4-curv}，$c_1(\cL)=\bigl[\frac{\ii}{2\pi}\omega\bigr]$；
若把 $\omega$ 的定向取反，$\cL$ 变负。半形式的平移符号依赖这一选择（D-4.7）。
\item \textbf{“行列式线丛”有两种。} 一种是 $\bar\partial$-族的 $\det$（本文 \src{369}），另一种是
Hodge 线丛 $\det H^0(\Sigma;\Omega^{1,0})^{\otimes}$。两者在 $\cM$ 上相差 $K_{\cM}$ 的幂，
混用会把 $h^\vee$ 的平移算错。
\item \textbf{$\cM$ 有奇点。} “$\cM_J$ 上的全纯截面”在奇点处要按 orbifold/分层理解；紧 Kähler 流形上
$H^0$ 有限维的定理不能直接用在奇点上，需要把 $\cM$ 当作射影簇（Narasimhan--Seshadri 保证它是）。
\end{enumerate}
\end{keypoint}

\begin{keypoint}
\textbf{一句话记住本章：} 三级台阶里，本文只显式做了第二级（极化），第一级（整性）用行列式线丛
“绕过”，第三级（半形式）完全没做、只在 §2 用 $\eta$ 不变式等价地实现。读 §3.1 时若把这三件事
当成一件，后面 §4 的 $k\to k+h^\vee$ 就无法解释了。
\end{keypoint}

\subsection{练习}

\begin{enumerate}[leftmargin=2.2em]
\item \textbf{（复算型）} 验证 \eqref{eq:ch4-preq} 与 $[\nabla_X,\nabla_Y]-\nabla_{[X,Y]}=-\ii\hbar\omega(X,Y)$
相互等价（提示：曲率 $F_\nabla=\nabla^2$，作用在截面上即对易子；比较两边的 2-形式）。

\item \textbf{（复算型）} 用 \eqref{eq:ch4-preq} 与陈类整性，推出 $\frac{1}{2\pi\hbar}[\omega]\in H^2(M;\ZZ)$，
并说明为什么这个条件\emph{只是}必要条件。

\item \textbf{（推导型）} 在 $\mathbb{CP}^1$ 上取 $\omega$ 的面积为 $A$。写出可预量子化的条件
（提示：$A\in\hbar\mathbb Z$ 的一半），并算出 $H^0(\mathbb{CP}^1,\cO(k))$ 的维数，
说明它等于 $k+1$ 而非 $k$ 的原因（提示：$\cO(k)$ 的截面对应 $k$ 次齐次多项式）[李思 p.307--308]。

\item \textbf{（推导型）} 证明：Kähler 极化下 $\nabla^{0,1}$ 平方为零 $\iff$ $F^{0,2}=0$ $\iff$
Kähler 形式闭。指出哪一步用到了 $(M,J,\omega)$ 的 Kähler 性。

\item \textbf{（推导型）} 在旗流形 $G/T$ 上验证 \eqref{eq:ch4-half}：用 $K_{G/T}=\cL_{-2\rho}$ 与
“$\cL_\lambda\otimes\cL_\mu=\cL_{\lambda+\mu}$” 得 $K^{1/2}=\cL_{-\rho}$，并说明它使最高权
$\lambda$ 的 Weyl 维数公式出现 $\lambda+\rho$。

\item \textbf{（判断题）} “本文没有讨论整性，所以 CS 量子化不需要整性条件。”这句话错在哪？
给出你在 D-4.2 中找到的两条理由。

\item \textbf{（判断题）} “半形式修正给出 $k\to k+h^\vee$，所以本文的半形式被作者略去了。”
这句话一半对一半错：对的是本文确实没做半形式；错的是把平移的\emph{来源}归给半形式。
依据本文，平移的正确来源是什么？（提示：§2 的 (2.16)(2.17)。）

\item \textbf{（判断题）} 若把 $\omega$ 换成 $-\omega$，哪一行的公式会变号？列出至少两处
（提示：\eqref{eq:ch4-preq} 与 \eqref{eq:ch4-curv}）。
\end{enumerate}

\subsection{本章边界}

\begin{boundary}
\begin{itemize}[leftmargin=1.4em]
\item \textbf{原文给出的}：$\cM_J$ 的全纯解释、行列式线丛 $\cL$、$k$ 次幂即预量子化线丛、
$H^0(\cM_J,\cL^{\otimes k})$ 为量子态空间及其有限维性 \src{368--369}。
\item \textbf{我们补出的}：D-4.1（Weil--Kostant 陈述与整性）、D-4.2（构造性绕过）、
D-4.3（Quillen 度规与曲率，陈述级）、D-4.4--D-4.5（极化三型与全纯截面）、D-4.6（Narasimhan--Seshadri）、
D-4.7（半形式与 $k+h^\vee$ 的两条证据）、D-4.8（$k$ 的三重身份）。
\item \textbf{仍是陈述级引用}：Quillen 度规曲率 \eqref{eq:ch4-curv}（Freed 1987 等，未逐字核）；
Narasimhan--Seshadri 的定理陈述（未核页码）；$K_{\cM_J}\simeq\cL^{-2h^\vee}$ \eqref{eq:ch4-kmod}。
\item \textbf{本文明确没做、我们也没有独立推出的}：由半形式机制\emph{独立}推出 $k\to k+h^\vee$。
本书的说法是“两条独立证据相符”，不是“我们证明了”。
\item \textbf{未核}：原文脚注 9 的完整内容（只核到首句与 $E_8$ 的例子）；河野 2.5 的半形式公式
（只核到章节位置）；Moore 3.4 的页码为页脚页，未逐句核。
\end{itemize}
\end{boundary}

\noindent\textbf{G4 覆盖清单}
\begin{table}[htbp]
\centering
\small
\setlength{\tabcolsep}{4pt}
\begin{tabular}{@{}lll@{}}
\toprule
缺口 & 本章位置 & 状态 \\
\midrule
G4.1 Weil--Kostant 陈述与本文的绕过 & D-4.1, D-4.2 & 已补（[我们]）\\
G4.2 行列式线丛、Quillen 度规、曲率 $\frac{\ii}{2\pi}\omega$ & D-4.2, D-4.3 & 已补（陈述级）\\
G4.3 极化定义、Kähler 极化、Bohr--Sommerfeld & D-4.4, D-4.5 & 已补（[我们]）\\
G4.4 量子化 = 协变常数截面、有限维 & D-4.5 & 已补（引 [李思 p.306--307]）\\
G4.5 半形式与 $k\to k+h^\vee$ & D-4.7 & 已补；本文未做，明说\\
G4.6 Narasimhan--Seshadri & D-4.6 & 已补（页码未核）\\
\bottomrule
\end{tabular}
\caption{G4.1--G4.6 覆盖清单。}
\end{table}

